\clearpage\fancyhead[L]{\small PALLAS.}

\begin{center}\textbf{PALLAS.}\par\vspace{4pt}\hrule\vspace{4pt}\fontsize{10}{12}\selectfont Göttingische gelehrte Anzeigen. Issue 32. pp. 305\,..\,308. 24 February 1810.\par\vspace{4pt}\hrule\vspace{4pt}\end{center}

\par\indent In 1808 \emph{Pallas} had been observed at only a few observatories, and only imperfectly. No revision of the elements had therefore been possible since 1807. At our observatory the planet’s extreme faintness made any observations with the mural quadrant impossible. The few observations made with the ring micrometer could make no claim to great accuracy, owing to particular circumstances. At the Seeberg observatory, Mr. \textsc{von Lindenau} did observe the planet with the transit instrument with his usual precision, but declinations could not be observed there either. The determinations made with the equatorial sector in Milan begin only on 22 August, so they could hardly be used for that opposition. Only through the Astronomical Yearbook for 1812 did a few complete observations by Mr. \textsc{David} in Prague around the time of opposition become known. Professor \textsc{Gauss} used their declinations together with the right ascensions determined at the Seeberg observatory to derive the opposition. Unfortunately, examination of the Prague observations showed poor agreement. For example, the different declinations yield results that differ by nearly one minute. Since there were no better declinations, the available ones had to be used as well as possible. The following result was thus obtained for the 1808 opposition, the fifth observed so far.

\begin{center}Opposition of Pallas, 26 July 1808.\\$21^{\text{u}}\;17^{\text{m}}\;32^{\text{s}}$ Göttingen meridian\par True longitude .\;.\;.\;.\;.\; $304^\circ\;2'\;59''7$\par True geocentric latitude\quad $37\;33\;54$ north\end{center}

\Needspace{26\baselineskip}\SourceLeaf{323}\par\indent In 1809 the observations of Pallas made at our observatory (see Gel. Anz. 1809, p. 1348) [26 Aug. 1809, p. 314 of this volume] seem to have been the earliest. From other observatories, Professor \textsc{Gauss} received a valuable series of Seeberg and Paris meridian observations through the kindness of Mr. \textsc{von Lindenau} and Mr. \textsc{Bouvard}. We give the Paris observations below, as they would otherwise become known only after a few years. Professor \textsc{Gauss} himself added the mean times, which Mr. \textsc{Bouvard} had not supplied.

\Needspace{18\baselineskip}\begin{center}\emph{Observations of Pallas at Paris.}\end{center}
\begin{center}\fontsize{9}{11}\selectfont\setlength{\tabcolsep}{5pt}\renewcommand{\arraystretch}{1.0}\begin{tabular}{lrrr}
1809 & Mean time & \shortstack{Apparent right\\ascension} & \shortstack{Apparent southern\\declination}\\
Sept. 12 & $12^{\text{u}}\;51^{\text{m}}\;39^{\text{s}}6$ & $4^\circ\;30'\;25''80$ & $4^\circ\;20'\;10''4$\\
15 & $12\;37\;41.7$ & $3\;57\;48.75$ & $5\;5\;16.4$\\
28 & $11\;36\;52.6$ & $1\;31\;50.25$ & $8\;20\;23.4$\\
29 & $11\;32\;11.1$ & $1\;20\;24.30$ & $8\;34\;57.3$\\
Oct. 2 & $11\;18\;7.1$ & $0\;46\;15.30$ & $9\;18\;22.5$\\
5 & $11\;4\;5.9$ & $0\;12\;35.25$ & $10\;0\;21.0$\\
6 & $10\;59\;25.5$ & $0\;1\;37.50$ & $10\;14\;6.9$\\
7 & $10\;54\;46.2$ & $359\;50\;44.55$ & $10\;27\;36.8$\\
8 & $10\;50\;7.1$ & $359\;39\;54.90$ & $10\;40\;1.5$\\
9 & $10\;45\;28.7$ & $359\;29\;16.50$ & $10\;54\;5.5$\\
13 & $10\;27\;2.2$ & $358\;48\;26.25$ & $11\;54\;59.5$\\
\end{tabular}\end{center}

\par\indent The numbers printed here agree exactly with those in Mr. \textsc{Bouvard}’s letter. Analysis of this determination has shown, however, that the right ascension for 12 September is about $1'\,15''$ of arc too large (the right ascensions were given in time only), that the declination for 8 October is 1 minute too small, and that the declination for 13 October is 10 minutes too large. These observations give the following determinations for the sixth opposition of Pallas.

\begin{center}Time of opposition, 22 Sept. 1809.\\$16^{\text{u}}\;10^{\text{m}}\;20^{\text{s}}$ Göttingen meridian\par True longitude of Pallas .\;.\; $359^\circ\;40'\;4''4$\par True geocentric latitude ,\quad $7\;22\;10.1$ south\end{center}

\par\indent The results Professor \textsc{Gauss} derived from these two oppositions in conjunction with the four earlier ones will be published on another occasion.

As we did last year, we also print here a list recently sent to us for publication. It gives some stellar occultations visible in 1810 that the astronomical ephemerides either leave out or give incorrectly. Mr. \textsc{Wisniewsky} calculated them in advance for the meridian and horizon of Berlin. --- ---
